\chapter{The REPL}

\section{What is the REPL}

The REPL (\textit{Read-Eval-Print Loop}) is the interactive mode of
\hayashi{}. Each line typed is evaluated immediately and the result is
displayed. It is the ideal environment for exploring data, testing
expressions, and prototyping analyses.

\begin{lstlisting}[style=inline, language=sh]
hay
\end{lstlisting}

\begin{lstlisting}[style=inline]
Hayashi 0.2.4  — Applied Econometrics Language
In honor of Fumio Hayashi. Type 'exit' or Ctrl-D to quit.

hay>
\end{lstlisting}

\section{Exiting the REPL}

\begin{center}
\begin{tabular}{ll}
\toprule
\textbf{Command} & \textbf{Effect} \\
\midrule
\texttt{exit}   & Exit the REPL \\
\texttt{Ctrl-D} & Exit the REPL (EOF) \\
\texttt{Ctrl-C} & Cancel the current line \\
\bottomrule
\end{tabular}
\end{center}

\section{Navigation and history}

The REPL uses \texttt{rustyline} and supports:

\begin{itemize}
  \item Up/down arrow: command history
  \item \texttt{Ctrl-R}: history search
  \item \texttt{Ctrl-A} / \texttt{Ctrl-E}: beginning/end of line
  \item \texttt{Ctrl-C}: cancel the current line
  \item \texttt{Ctrl-D}: exit the REPL (equivalent to \texttt{exit})
\end{itemize}

History is persisted across sessions in \texttt{\$HOME/.hay\_history}.

\section{Expressions in the REPL}

Any valid expression can be typed directly:

\begin{lstlisting}[caption={Basic interactive session}]
hay> 2 + 2
4

hay> let x = 10
hay> x * 3
30

hay> let s = "hello, world"
hay> s
"hello, world"
\end{lstlisting}

\section{Running scripts}

\texttt{.hay} scripts are text files with one statement per line.
To run one:

\begin{lstlisting}[style=inline, language=sh]
hay analysis.hay
\end{lstlisting}

Output goes to \texttt{stdout}. Errors go to \texttt{stderr}.

\section{Error format}

\hayashi{} displays errors with visual context:

\begin{lstlisting}[style=inline]
error: line 3: undefined variable 'y' — did you mean 'x'?
   3 | display y + 1
     | ^^^^^^^^^^^^^
\end{lstlisting}

The code snippet, line number, and suggestion (when available) make
correction straightforward without having to count lines manually.
